% !PRINTED 125
\weberpage{125}
%<<CONT
\begin{center}
\S.\,30.\\
\gesperrt{Der zweite Sylow'sche Satz.}
\end{center}

Es sei nun, wie oben, $P$ eine Gruppe vom Grade $n$ und $p^{\varkappa}$ die höchste Potenz der Primzahl $p$, die in $n$ aufgeht; $Q$ sei ein Theiler von $P$ vom Grade $p^{\varkappa}$, der nach dem Satze I, \S.\,29 existirt.

Der Inbegriff aller Elemente $c$ von $P$, die der Bedingung genügen:
\weq{1}{c^{-1}Qc = Q,}
unter denen gewiss alle Elemente von $Q$ enthalten sind, bildet eine Gruppe; denn aus
\[ Qc = cQ,\quad Qc' = c'Q \]
folgt, dass auch $Qcc' = cc'Q$ ist. Diese Elemente $c$ können wir \gesperrt{die mit der Gruppe $Q$ vertauschbaren Elemente von $P$} nennen. Bezeichnen wir die Gruppe der $c$ mit $R$, so ist $R$ ein Theiler von $P$, und $Q$ ein Theiler von $R$, und zwar ein Normaltheiler. In besonderen Fällen kann $R$ sowohl mit $Q$, als mit $P$ identisch sein. Ist aber $b$ irgend ein nicht in $R$ enthaltenes Element von $P$, so giebt es unter den Elementen der Gruppe $b^{-1}Qb$ gewiss wenigstens eines, das nicht in $Q$ enthalten ist.

Bezeichnen wir mit $r$ den Index von $Q$ in Bezug auf $R$, mit $j$ den Index von $R$ in Bezug auf $P$, so ist
\weq{2}{n = p^{\varkappa}rj,}
und $r$ und $j$ sind durch $p$ nicht theilbar.

Der Satz, den wir zu beweisen haben, lautet:

\gesperrt{II. Der Index $j$ von $R$ ist von der Form $pk + 1$, worin $k$ irgend eine ganze Zahl ist $[j \equiv 1\ (\mathrm{mod}\ p)]$; es giebt $j$ und nicht mehr verschiedene Theiler von $P$ vom Grade $p^{\varkappa}$, die alle mit einander conjugirt sind. Jeder Theiler von $P$, dessen Grad eine Potenz von $p$ ist, ist in einer dieser conjugirten Gruppen enthalten.}

Ist $c$ ein Element aus $R$ vom Grade $\gamma$, und $s$ der kleinste positive Exponent, für den $c^s$ in $Q$ enthalten ist, so muss jeder andere Exponent $t$, für den $c^t$ zu $Q$ gehört, durch $s$ theilbar
%CONT>>
